\section{Formulare}
\label{modul:formulare}
Das Modul \module{Formulare} befindet sich noch in der Entwicklung. Es soll die Möglichkeit bieten, Formulare aus dem \pkg{hyperref} in PDF-Dokumenten komfortabel einzubinden. Dabei sollen auch die Möglichkeit erhalten bleiben, dass die gesetzten Dokumente in ausgedruckter Form für den Unterricht nutzbar sind.

Durch Einbindung des Moduls wird automatisch eine \texttt{Form}-Umgebung um das komplette Dokument erstellt. Die Übergabe von Parametern an diese Umgebung muss er noch implementiert werden.

\subsection{Befehle}
\begin{commands}
    \command{feldLinFormular}[\oarg{Abstand}\marg{Anzahl}\oarg{Textfeldargumente}]
		erstellt ein Textfeld für das Formular, dass mit der gegebenen Anzahl an Linien für den Ausdruck versehen ist. Der Abstand der Linien zueinander kann angegeben werden, er beträgt standardmäßig $0,7\unit{cm}$. Dem Befehl können auch Argumente mitgegeben werden, die an das intern genutzte \texttt{TextField} übergeben werden.

		Bei dem Textfeld ist die Schriftgröße 18pt gesetzt, die ungefähr passend zu den Linienabstand von $0,7\unit{cm}$ ist. Außerdem wurde die Hintergrundfarbe auf leer gesetzt, damit die Linien zu erkennen sind.
\begin{example}
  \feldLinFormular[0.5cm]{3}
\end{example}
\end{commands}